\chapter{Correspondence with the papers}
\label{chap:correspondence}

The endpoint of this account is the finite sticky estimate, WZ2
Theorem 5.2. In particular, WZ2's earlier reduction from general
Assouad configurations to all-scale CWA, and its final dimension
arguments, are outside the present scope. The other papers enter
through the structural, projection, incidence, and multilinear estimates
used to prove that endpoint.

\section{Manuscripts and numbering}

We use the following freely available versions. The numbering in a
Lean declaration may refer to an earlier manuscript or to the published
article; mathematical content takes precedence over that name.

\begin{center}
\begin{tabular}{lll}
Reference & arXiv version & Date \\
\hline
WZ1 & 2210.09581v2 & October 2025 \\
WZ2 & 2401.12337v2 & April 2025 \\
OSW25 & 2511.21656v1 & November 2025 \\
OS & 2106.03338v3 & September 2023 \\
OSW & 2209.00348v1 & September 2022 \\
PYZ & 2207.02259v3 & October 2024 \\
CV & 1205.6371v1 & May 2012
\end{tabular}
\end{center}

In OS, the regular-case result called Theorem 6.1 in the repository
is Theorem 6.2 in this manuscript. The combining result called
Proposition 7.3 is Proposition 7.5 here. The uniformization lemma is
Lemma 7.3 and the multiscale decomposition is Proposition 8.1.
The radial bootstrap cited in WZ1 as OSW Lemma 2.9 is Lemma 2.8
in the available arXiv manuscript, dated September 2022.
CV's multilinear theorem appears as WZ1 Theorem 4.2 in the version
used here; some Lean comments call it Theorem 10.

\section{Transfer from WZ1 to WZ2}
\label{sec:wz1-wz2-transfer}

The structural part of WZ2 reuses the WZ1 argument after changing its
extremal class. The following correspondences concern mathematical
proofs, rather than an identification of Lean declarations or a use
of an unchanged directional theorem on a CWA family.

\begin{center}
\begin{tabular}{p{.23\textwidth}p{.22\textwidth}p{.44\textwidth}}
WZ1 source & WZ2 result & Transfer needed \\
\hline
Proposition 3.2 & Proposition 6.2 & Rebuild all-scale CWA covers and
retain the family-cardinality factors in both multiplicity bounds. \\
Proposition 4.1 & Proposition 6.3 & Carry planiness and grain counts
through the balanced coarse and fine CWA configurations. \\
Proposition 5.1 & Proposition 6.4 & Apply the projection dichotomy,
AD exponent gap, and trapezoid smoothing to the retained grains. \\
Lemma 6.3 & Lemma 6.6 & Replace the final directional count by a
convex-prism containment count; the popular-slab preparation remains. \\
Proposition 6.1 & Proposition 6.5 & Normalize the derivative while
transporting the same shading, plane map, and AD bounds. \\
Corollary 6.5 & Lemma 6.7 & Retain the small projected-area conclusion
needed for the WZ2 contradiction. \\
Proposition 7.1 & Propositions 7.1, 6.8 & Use parameter extraction,
block replication, and finite scale iteration for the WZ2 lower bound.
\end{tabular}
\end{center}

The supplied Theorem 5.2 dependency atlas is useful for locating these
transfers and the intermediate proof operations. Its published-paper
page numbers are separate from the manuscript versions listed above.
Its Lean blocks are uncompiled interface sketches. Declaration links
here refer to the repository's audited names, and the added mathematical
steps do not acquire completion marks from those sketches.

Several qualifications matter when using the route. The normalized
Wolff count changes under tube deletion
(Lemma~\ref{lem:cwa-refinement}); each surviving containment fibre
has its own loss. The Section 7 replication used in Lean retains a
local cardinality deficit (Remark~\ref{rem:replication-deficit}).
The final parameter-Frostman theorem also requires explicit indexed
parameter counts (Remark~\ref{rem:parameter-frostman}). These are
part of the passage between the displayed mathematical interfaces.

The atlas stops at the external multilinear, radial, cinematic, and
entropy inputs. This blueprint also opens the supporting arguments
implemented in the repository. The WZ1 continuum statements and
self-similar CWA remain contextual statements; the finite endpoint
does not gain them as additional dependencies.

\section{Coverage and formal scope}

The geometric definitions from WZ2 Section 2, the sticky endpoint,
all the numbered propositions and lemmas of WZ2 Sections 6--7,
and the WZ1 structural and projection arguments are recorded in the
preceding chapters. OSW25's seventeen numbered results are included.
The OS chapter follows the discrete incidence and multiscale argument,
and the radial chapter records the bootstrap and its WZ1 applications.
For PYZ, we include the cinematic estimate, the curve-generated
specialization, and the geometric and counting statements in
Sections 3--5. The multilinear chapter includes CV's numbered
ingredients, sometimes grouping the original and mollified versions.
Continuum projection consequences not used in the sticky proof are
not developed separately.

There are three different kinds of correspondence. Some declarations
prove the displayed statement outright, as with the final sticky
theorem. Others prove precisely the normalized case used in the
argument, such as the three-dimensional multilinear estimate and the
$L^{3/2}$ cinematic input. Still others define an interface or prove
one implication in a larger argument. In the last case, the link is a
way to inspect that interface, not a claim that its entire paper
statement has been proved by the linked declaration alone.

The general PYZ Theorem 1.7 and the arbitrary-dimensional CV theorem
are useful mathematical context, but the completed interfaces used
by this project have the narrower scopes stated in their chapters.
Similarly, the normalized strong ring theorem and the separated
projection result retain the parameter ranges and endpoint alternatives
actually supplied by Lean. An unmarked node may have substantial
formal work behind it; absence of a completion mark is not evidence
that the corresponding mathematics is absent from the repository.

The compiled dependency audit provides a separate check of usage. For
example, the sticky proof reaches the cone version of the Marstrand
estimate and the uniform expansion lemma. Its Frostman-to-Katz--Tao
step uses the weighted-cardinality theorem, while its multilinear
input passes through the broad-set estimate for paper tube shadings.
These choices need not pass through every other proved formulation
linked in the mathematical account. The generated comparison table
records each declaration separately and supplies a source-level
dependency path for every reached version.

\section{Changes that affect a statement}

The remarks at the relevant statements give the mathematical details.
The following list locates the main issues without treating every
normalization as an error in a paper.

\begin{itemize}
\item The sticky endpoint uses nonempty families and aggregate shading
density. Its public hypothesis is the literal nearby-scale CWA
condition, including full containment fibres and unit rescaling.
See Remarks~\ref{rem:endpoint}, \ref{rem:aggregate-density},
\ref{rem:full-fibres}, and \ref{rem:nearby-scale}.
\item A balanced cover carries equal fine-union volume in each retained
coarse cube. This is stronger than pointwise shading compatibility or
uniform cardinality of containment fibres; see
Remark~\ref{rem:balanced-cover}.
\item The plane map is evaluated on the shaded union, and its normal
transforms as a covector. The AD condition is an upper covering bound,
not two-sided Ahlfors--David regularity; see
Remark~\ref{rem:upper-ad} and the local-grain discussion.
\item The small twisted image has exponent $\sigma-\eps$ in the
formal argument, whereas WZ2 Lemma 6.7 prints $\sigma+\eps$.
The minus sign is enough for the contradiction. The same chapter
records the reversed ratio in the printed Katz--Tao condition and
the scale conventions in the uniformization and replication lemmas.
The replication deficit, the distinction between the parameter-block
scale and its rescaled mesh, and the finite iteration budget are made
explicit in Remarks~\ref{rem:replication-deficit},
\ref{rem:projection-scales}, and \ref{rem:projection-iteration}.
\item In the dot-product dichotomy, it is the coefficient set that
concentrates near a line through zero. The endpoint set concentrates
near an affine line. The closed separated result yields concentration
in at least one endpoint class. The distinction is essential when the
classes are independent.
\item Passing from a measure-valued thin-tube condition to a statement
for every retained discrete point changes the exceptional fraction
and the counting constant. These losses are kept in the radial chapter.
\item The cinematic estimate used here fixes an absolute $C^2$ bound
before choosing the small-scale threshold. Pairwise cinematic
conditions alone do not impose such a bound on a singleton family.
The original WZ1 projection-norm upper bound also needs a bound on
the slope offset $f(0)$; see Remark~\ref{rem:wz1-slope-offset}.
\item OSW25's Marstrand inequality is a lower bound for average
projection length by inverse energy. Its Ruzsa Corollary 2.4 has a
different denominator in the general two-set Lean formulation; the
equal-set specialization agrees.
\item The CV degree comparison needs a positive degree bound, and
the ellipsoid colouring keeps a quantitative separation margin.
See Remarks~\ref{rem:cv-polynomial} and \ref{rem:cv-colouring}.
\end{itemize}

These distinctions preserve the logic of the argument while making
clear exactly what each formal declaration says. The full theorem is
the declaration
\nolinkurl{Kakeya.Assouad.PureWZ2Theorem5_2Unconditional}; its statement
has no external proof inputs.
